\section{离散扩散模型：蛋白质序列的加噪与去噪}

\subsection{离散数据的加噪策略}

对于离散序列数据（如自然语言或蛋白质序列），Amini 提出了两种实现"噪声化"的策略：

\textbf{策略一：Masking（遮蔽）}

从干净序列出发，在每一步随机选择一个或多个位置，将对应的 token 替换为特殊的 \texttt{[MASK]} 标记：

\begin{table}[H]
\centering
\caption{Masking 策略示意（以自然语言"This class is fun"为例）}
\begin{tabular}{@{}cl@{}}
\toprule
\textbf{步骤} & \textbf{序列状态} \\
\midrule
$t=0$ & This class is fun \\
$t=1$ & This {[}MASK{]} is fun \\
$t=2$ & {[}MASK{]} {[}MASK{]} is fun \\
$t=3$ & {[}MASK{]} {[}MASK{]} {[}MASK{]} fun \\
$t=4$ & {[}MASK{]} {[}MASK{]} {[}MASK{]} {[}MASK{]} \\
\bottomrule
\end{tabular}
\end{table}

\textbf{策略二：Mutation（突变/替换）}

不使用特殊标记，而是在每一步以一定概率将某些位置的 token 随机替换为词汇表中的其他 token：

\begin{table}[H]
\centering
\caption{Mutation 策略示意}
\begin{tabular}{@{}cl@{}}
\toprule
\textbf{步骤} & \textbf{序列状态} \\
\midrule
$t=0$ & This class is fun \\
$t=1$ & This \textit{apple} is fun \\
$t=2$ & \textit{Rock} \textit{apple} is fun \\
$t=3$ & \textit{Rock} \textit{apple} \textit{my} fun \\
$t=T$ & （完全随机的 token 序列） \\
\bottomrule
\end{tabular}
\end{table}

\begin{warningbox}{Mutation 比 Masking 更难}
Mutation 策略对去噪模型的要求更高，因为模型无法通过特殊标记识别哪些位置被修改了——它需要自行判断哪些 token 是"错误的"并预测正确值。这在某种意义上更接近真实的进化突变过程，但也大大增加了训练难度。在 EvoDiff 中，两种策略都被采用和测试。
\end{warningbox}

\subsection{离散扩散的形式化}

对于蛋白质序列 $\mathbf{s} = (s_1, s_2, \ldots, s_L)$，其中 $s_i \in \{A_1, A_2, \ldots, A_{20}\}$ 表示 20 种氨基酸之一：

\textbf{前向过程}：在每一步 $t$，随机选择一个子集的位置 $\mathcal{M}_t$，将这些位置上的氨基酸替换为 \texttt{[MASK]}：

$$s_i^{(t)} = \begin{cases} \texttt{[MASK]} & \text{if } i \in \mathcal{M}_1 \cup \mathcal{M}_2 \cup \cdots \cup \mathcal{M}_t \\ s_i^{(0)} & \text{otherwise} \end{cases}$$

\textbf{反向过程}：训练神经网络 $f_\theta$，在每一步将一些被 mask 的位置恢复为正确的氨基酸：

$$f_\theta(\mathbf{s}^{(t)}, t) \rightarrow \mathbf{s}^{(t-1)}$$

生成时，从全 mask 序列 $\mathbf{s}^{(T)} = (\texttt{[MASK]}, \texttt{[MASK]}, \ldots, \texttt{[MASK]})$ 出发，逐步去噪，最终得到一条完整的蛋白质序列。

\subsection{离散扩散与语言模型的关系}

Amini 指出一个深刻的理论联系：\textbf{离散 Masking Diffusion 实际上是多种语言建模方案的泛化}：

\begin{table}[H]
\centering
\caption{离散扩散与经典语言建模方案的对比}
\begin{tabular}{@{}lll@{}}
\toprule
\textbf{方法} & \textbf{解码顺序} & \textbf{每步预测} \\
\midrule
Next Token Prediction (LLM) & 固定的从左到右 & 单个下一 token \\
Masked Language Model (BERT) & 全序列可见 & 一步填充所有 mask \\
Discrete Diffusion & 随机顺序，逐步 & 每步填充部分 mask \\
\bottomrule
\end{tabular}
\end{table}

\begin{importantbox}{离散扩散的泛化能力}
离散扩散模型学习的是\textbf{所有可能的解码顺序}和\textbf{所有可能的 masking 子集}上的去噪任务。这使得它同时覆盖了：
\begin{itemize}
    \item 自回归模型的"从左到右"解码（一种特殊的 masking 顺序）
    \item Masked Language Model 的"一步去噪"（一种特殊的步数设置）
\end{itemize}
因此，离散扩散模型是一个更加通用的框架，具有更强的灵活性和表达能力。
\end{importantbox}

\subsection{本章小结}

将 Diffusion Model 从连续域迁移到离散域的关键在于重新定义"噪声"——通过 Masking 或 Mutation 对序列 token 进行逐步破坏。离散扩散框架在数学上统一了自回归模型和 Masked Language Model，为蛋白质序列建模提供了强大且灵活的工具。


